\section{Scope, principal results, and relation to the source material}
\label{sec:scope}

This report continues the polylogarithm manuscript and the active incoming
research of the \repo\ project. It is based on the fixed repository commit
\begin{center}\small\texttt{\baselinecommit}\end{center}
The consolidated manuscript and the newer
continuation reports must be read together: several questions still visible
in older fragments have already been settled in later reports. The source
comparison, exact paths, and proposed integration actions are supplied in
the accompanying claim ledger and provenance manifest.

The results here fall into four connected developments. The first is a
translation calculus for the Hurwitz order jets that generate balanced
Stieltjes antiderivatives. The second resolves an explicit joint-limit
conjecture for Gaussian polylogarithmic Euler remainders. The third studies
the exact signs and critical sums of the Lerch deformation velocities
already identified in the latest continuation. The fourth gives an
algebraic reduction preserving depth in the formal Cayley quotient.
Each development has its own hypotheses; none is used as a substitute for
an unproved arithmetic relation among numerical periods.

\subsection{Results at a glance}

\begin{enumerate}
\item \textbf{Exact translated Hurwitz pairings.}
Theorem~\ref{thm:translation} expresses the bilinear periodic translation
of two Hurwitz functions as a Gamma--trigonometric kernel times a
symmetric Hurwitz difference. Its polylogarithmic form permits every
order derivative. Theorem~\ref{thm:cusp} identifies the complete
logarithmic endpoint polynomial and an ordinary convergent even power
series remainder at every derivative order.

\item \textbf{A sharp Gamma translation inequality.}
The squared $L^2$ change in $\log\Gamma$ under a periodic translation is
strictly concave on $(0,1)$ and has its unique maximum at $1/2$
(Theorem~\ref{thm:gamma-max}). Its three exact representations yield
an accelerated odd-zeta series for $\gamma_1$ with an explicit geometric
tail (Corollary~\ref{cor:gamma1-fast}). Higher derivatives connect this
energy to polygamma functions and Hurwitz order derivatives; explicit
primitives follow from the same kernel.

\item \textbf{Proof of the near-critical Euler conjecture.}
For outer order $a=1-\varepsilon$, fixed inner order $b>0$, and
$\lambda=\log(1/\varepsilon)$, the unique reversal index satisfies
\[
 \nu(1-\varepsilon,b)\sim
 \frac{\lambda^{2b+2}}
 {\pi b^2\Gamma(b+1)^2\zeta(b+1)^2\varepsilon^2}.
\]
Theorem~\ref{eutrans:thm:crossing} proves the source conjecture
\texttt{research:conj:crossing}, including a first inverse-logarithmic
correction uniformly for $b$ in compact positive intervals. The proof
retains the exponentially small seed and an $\varepsilon$ factor in
the algebraic remainder. A rescaled profile is
$x^{-1/2}-x^{-1}$; the later maximum occurs at an index asymptotic
to $4\nu$.

\item \textbf{Optimal sign blocks and all-order critical finite parts.}
For the Lerch velocities $v_n(A)$, let $\theta(A)$ be the phase of
the dominant inverse singularity. The least eventual block length
containing both signs is exactly $\lceil\pi/\theta\rceil+1$
(Theorem~\ref{newdiag:thm:sharp-block}). At $A=2$ it is four.
Theorem~\ref{newdiag:thm:finiteparts} computes the finite part of every
sharp sum $\sum_{n\leq N}n^pv_n(A)\tau^n$. Both conjugate singularities
are required for the divergent subtractions once $p\geq2$.

\item \textbf{A depth-preserving formal Cayley reduction.}
An oriented choice of shuffle generators yields a multiplicative
retraction whose kernel is the full Cayley ideal and whose output
does not increase depth. Its generator counts give the exact depth
filtration of the quotient. This is a structural algebraic result;
the mixed $S_6$ period identity remains conjectural.
\end{enumerate}

\subsection{Classical antecedents and what is claimed here}

Hurwitz's Fourier expansion, the parameter recurrence, and the
order-zero Gamma evaluation are classical \cite{DLMFHurwitz}.
Espinosa and Moll developed product integrals and antiderivatives involving
Hurwitz zeta \cite{EM1,EM2}, and a generalized polygamma framework
\cite{EMPoly}. These works supply the appropriate context for the first
two sections. We give a direct analytic derivation of the translated
identities and their sharp consequences; the underlying unshifted product
integral is explicitly credited. The standard polylogarithm conventions
are those summarized in \cite{DLMFPolylog}.

The Euler section uses the exact density and the global one-crossing
theorem already proved in the newest applicable source \cite{RepoEuler}.
Its new work is the uniform estimate in the singular joint limit, the
corrected root law, and the profile consequences. The Lerch sections use
the selected-sheet continuation and additive coefficient asymptotic from
\cite{RepoLerch}; they answer two further questions posed there. Lagrange
inversion and singularity transfer are standard tools
\cite{Gessel,FlajoletOdlyzko}. Earlier Lambert polynomial and branch
expansions provide further context \cite{KaluginJeffrey,Cohen}.

Throughout, ``new'' means established here relative to the checked
repository baseline. This source audit is not a claim of worldwide
priority for every identity or special case. The classical tools are
identified, inherited repository theorems are marked, and new assertions
are supplied with proofs. The unresolved $S_6$ and current $S_8$ candidates
retain their conjectural status even when numerical residuals are tiny.

\subsection{Conventions and logical boundaries}

The Laurent convention is
$\zeta(1+u,a)=u^{-1}+\sum_{n\geq0}(-1)^n\gamma_n(a)u^n/n!$;
in particular $\gamma_0(a)=-\psi(a)$. Superscripts on $\zeta$ denote
derivatives with respect to its order, while superscripts on
$\gamma_n$ and $\psi$ denote derivatives with respect to their real
argument. Rising factorials are $(z)_m$ and falling factorials are
$(z)_{\underline m}$. Endpoint values of periodic $L^2$ functions are
understood almost everywhere. Branches of complex powers and square
roots are fixed where first used.

The analytic arguments, exact rational algebra, and floating-point
diagnostics play different roles. Proofs justify every limiting operation
that enters a theorem. Exact finite calculations check implementations
and stated ranks in specified spaces. Floating-point calculations test
formulas independently and illustrate asymptotic accuracy; they are not
interval certificates. The package does not claim proof-assistant
formalization or arithmetic independence from numerical experiments.
